\documentclass[11pt]{amsart}
\usepackage{calc,amssymb,amsmath,ulem,stmaryrd,fullpage}
\usepackage{enumerate}
\usepackage{alltt}
\RequirePackage[dvipsnames,usenames]{color}
\let\mffont=\sf
\normalem
%\input{kmacros3.sty}
\input{xy}
\xyoption{all}
\usepackage{hyperref}
\usepackage{graphicx}
\usepackage[all,cmtip]{xy}
\usepackage{verbatim}
\usepackage[left=0.95in,top=0.95in,right=0.95in,bottom=0.95in]{geometry}
\newcommand{\tauCohomology}{T}
\newcommand{\FCohomology}{S}
\usepackage{mathrsfs}


\theoremstyle{theorem}
%\newtheorem*{mainthm*}{Main Theorem}

\newcommand{\note}[1]{\marginpar{\sffamily\tiny #1}}

\def\todo#1{\textcolor{Mahogany}%
{\sffamily\footnotesize\newline{\color{Mahogany}\fbox{\parbox{\textwidth-15pt}{#1}}}\newline}}

\renewcommand{\O}{\mathcal O}
\newcommand{\ie}{{\itshape i.e.} }
\newcommand{\roundup}[1]{\lceil #1 \rceil}
\newcommand{\rounddown}[1]{\lfloor #1 \rfloor}
\renewcommand{\to}[1][]{\xrightarrow{\ #1\ }}
\newcommand{\onto}[1][]{\protect{\xrightarrow{\ #1\ }\hspace{-0.8em}\rightarrow}}
\newcommand{\into}[1][]{\lhook \joinrel \xrightarrow{\ #1\ }}
%\newcommand{DBDual}[1]{{\underline{\omega}_{#1}}}
\newcommand{\DBDual}[1]{{{\uuline \omega}{}^{\mydot}_{#1}}}
\newcommand{\Tt}{{\mathfrak{T}}}
%\newcommand{\bn}{{\mathfrak{n}} }
\newcommand{\sol}[1]{ \vskip 6pt{\color{blue} {\bf Solution:} #1} \vskip 6pt}

\newcommand{\bR}{{\mathbb{R}}}
\newcommand{\va}{{\vec{a}}}
\newcommand{\vb}{{\vec{b}}}
\newcommand{\vzero}{{\vec{0}}}
\newcommand{\vi}{{\vec{i}}}
\newcommand{\vj}{{\vec{j}}}
\newcommand{\vk}{{\vec{k}}}
\newcommand{\vh}{{\vec{h}}}
\newcommand{\vr}{{\vec{r}}}
\newcommand{\vu}{{\vec{u}}}
\newcommand{\vv}{{\vec{v}}}
\newcommand{\vw}{{\vec{w}}}
\newcommand{\vx}{{\vec{x}}}
\newcommand{\vy}{{\vec{y}}}
\newcommand{\vz}{{\vec{z}}}
\newcommand{\vT}{{\vec{T}}}
\newcommand{\vN}{{\vec{N}}}
\newcommand{\vF}{{\vec{F}}}
\newcommand{\vG}{{\vec{G}}}
\newcommand{\bF}{\mathbb{F}}
\newcommand{\bQ}{\mathbb{Q}}
\newcommand{\bZ}{\mathbb{Z}}
\newcommand{\bA}{\mathbb{A}}
%\newcommand{\cF}{\mathcal{F}}
%\newcommand{\cG}{\mathcal{G}}
%\newcommand{\cH}{\mathcal{H}}
\newcommand{\image}{\mathrm{im}\,}
\newcommand{\coker}{\mathrm{coker}\,}
\newcommand{\Hom}{\mathrm{Hom}\,}
\newcommand{\Supp}{\mathrm{Supp}\,}
\newcommand{\Ann}{\mathrm{Ann}\,}
\newcommand{\sH}{\mathscr{H}\,}
\newcommand{\cI}{\mathcal{I}}
\newcommand{\cJ}{\mathcal{J}}
\newcommand{\cE}{\mathcal{E}}
\newcommand{\cF}{\mathcal{F}}
\newcommand{\cG}{\mathcal{G}}
\newcommand{\cH}{\mathcal{H}}
\newcommand{\cO}{\mathcal{O}}
\newcommand{\cM}{\mathcal{M}}
\newcommand{\cN}{\mathcal{N}}
\newcommand{\cK}{\mathcal{K}}
\newcommand{\Div}{\mathrm{Div}}

\begin{document}
\title{Worksheet \#3 -- Math 6140\\Spring 2019}
\author{Due Monday, February 18th}

\maketitle
You may work in groups of up to 3.  Only one worksheet needs to be turned in per group.
\vskip 9pt
\noindent
We begin by recording some definitions (we'll talk about these soon).
\vskip 9pt
\noindent
{\bf Definition.}
We work on a normal quasi-projective variety $X$ and we fix $\cK(X)$ to be the constant sheaf made out of the function field $K(X)$ of $X$.
Suppose $D$ is a Weil divisor.  Set $\O_X(D) \subseteq \cK(X)$ denote the sheaf of $\O_X$-modules with the property that
\[
\Gamma(U, \O_X(D)) = \{f \in K(X) \;|\; \Div_U(f) + D \geq 0 \}.
\]

A divisor $D$ is called a \emph{Cartier divisor} if $\O_X(D)$ is an invertible sheaf.  In other words, for each $p \in X$, we have isomorphisms of stalks $\O_X(D)_p = \O_{X,p}$.
\vskip 12pt
\noindent
{\bf 1.}  Show that every principal divisor on a non-singular variety is Cartier.  
\vskip 180pt
\noindent
{\bf 2.}  Suppose that $D$ is a Cartier divisor and $E$ is a Weil divisor.  Prove that $$\O_X(D) \otimes_{\O_X} \O_X(E) \cong \O_X(D) \cdot \O_X(E) = \O_X(D+E) \subseteq \cK(X).$$
\newpage
\noindent
{\bf 3.}  Consider $X = Z(xy-z^2) \subseteq \bA^3$.  Verify that the ideal $\langle x, z \rangle \subseteq k[x,y,z]/\langle xy-z^2 \rangle = k[X]$ is prime of height one and so determines a prime divisor $E$.  Show that $E$ is not Cartier but that $2E$ is.
\vskip 200pt
Suppose that $\pi : Y \to X$ is a finite surjective morphism of normal affine varieties (so that $\pi^* : k[X] \to k[Y]$ makes $k[Y]$ into a finite $k[X]$-modules).  For every prime divisor $E$ on $X$, with corresponding prime ideal $I_E$ and discrete valuation $v_E$, there exist finitely many prime divisors $F_1, \dots, F_n \subseteq Y$ that map onto $E$.  Suppose that $e \in \O_{X, E}$ is an element with $v_E(e) = 1$.  Write $a_i = v_{F_i}(e)$.  We write
\[
\pi^* E = \sum_i a_i F_i.
\]
For more general divisors $\sum b_j E_j$, we extend by linearity.  
\vskip 12pt
\noindent
{\bf 4.}  Consider the map $\pi : \bA^1 \to \bA^1$ defined by $\pi(t) = t^2$.  Compute $\pi^* (2 \cdot Z(x) - 3 \cdot Z(x-1) + 1 \cdot Z(x+1))$.
\newpage
\noindent
{\bf 5.}  Suppose that $\pi : Y \to X$ is a finite surjective morphism of normal affine varieties and that $f \in K(X)$.  Show that 
\[
\pi^* \Div_X(f) = \Div_Y(f).
\]
\vskip 170pt
Suppose that $\pi : Y \to X$ is a surjective morphism of normal varieties.  Suppose that $D$ is a Cartier divisor on $X$.  For every sufficiently small open chart $U$ on $X$, since $\O_X(D)$ is locally free, we can write $f_i \cdot \O_{U_i} = \O_X(D)|_{U_i}$ for some $f_i \in K(X)$.  We define $\pi^* D$ to be the divisor $\Div_{\pi^{-1}(U_i)}(f_i)$ on $\pi^{-1}(U_i)$.
\vskip 12pt
\noindent
{\bf 6.}   Suppose that $\pi : Y \to X$ is a surjective morphism of normal affine varieties and $D$ is a Cartier divisor on $X$.  Prove that the two definitions of $\pi^* D$ agree.
\vskip 200pt
\noindent
{\bf 7.}  Consider $\pi : Y \to X = \bA^2$ to be the blowup of the origin.  Let $D = \Div_X( (y^2- x^3)/y)$.  Write $\pi^* D$ as a sum of prime divisors.  
\end{document}

